\documentclass[11pt,a4paper]{report}
\usepackage[utf8]{inputenc}
\usepackage[spanish]{babel}
\usepackage{graphicx}
\usepackage{makeidx}
\usepackage[pdftex,
			pagebackref=true,
            colorlinks=true,
            linkcolor=blue!70!black,
            urlcolor=blue!70!black,
            citecolor=blue!70!black,
            unicode
			]{hyperref}
\usepackage[margin=3cm]{geometry}
\usepackage{makeidx}
\usepackage{multicol}
\usepackage{multirow}
\usepackage[nottoc]{tocbibind}
\usepackage{longtable}
\usepackage[x11names,table]{xcolor}
\usepackage{amsmath,amssymb,amsfonts,latexsym,cancel}
\usepackage{amsthm}
\usepackage{color}
\usepackage{algorithmic}
\usepackage[chapter]{algorithm}
\usepackage{listings}
\usepackage{caption}
\usepackage{subcaption}
\usepackage{float}
\usepackage{forloop}
\usepackage{tikz}
\usepackage{xstring}
\usepackage{fancyhdr}
\usepackage{lipsum}

\pagestyle{fancy}
\lhead{\slshape\nouppercase{\rightmark}}
\rhead{\thepage}

%\renewcommand{\thesubfigure}{\arabic{subfigure}}

\lstset{ %
  xleftmargin=40pt,
  backgroundcolor=\color{white},   % choose the background color; you must add \usepackage{color} or \usepackage{xcolor}
  basicstyle=\footnotesize,        % the size of the fonts that are used for the code
  breakatwhitespace=false,         % sets if automatic breaks should only happen at whitespace
  breaklines=true,                 % sets automatic line breaking
  captionpos=b,                    % sets the caption-position to bottom
  commentstyle=\color{green},    % comment style
  deletekeywords={...},            % if you want to delete keywords from the given language
  escapeinside={\%*}{*)},          % if you want to add LaTeX within your code
  extendedchars=true,              % lets you use non-ASCII characters; for 8-bits encodings only, does not work with UTF-8
  frame=none,                    % adds a frame around the code
  keepspaces=true,                 % keeps spaces in text, useful for keeping indentation of code (possibly needs columns=flexible)
  keywordstyle=\color{blue},       % keyword style
  language=C++,                 % the language of the code
  morekeywords={*,...},            % if you want to add more keywords to the set
  numbers=left,                    % where to put the line-numbers; possible values are (none, left, right)
  numbersep=10pt,                   % how far the line-numbers are from the code
  numberstyle=\color{gray}, % the style that is used for the line-numbers
  rulecolor=\color{black},         % if not set, the frame-color may be changed on line-breaks within not-black text (e.g. comments (green here))
  showspaces=false,                % show spaces everywhere adding particular underscores; it overrides 'showstringspaces'
  showstringspaces=false,          % underline spaces within strings only
  showtabs=false,                  % show tabs within strings adding particular underscores
  stepnumber=1,                    % the step between two line-numbers. If it's 1, each line will be numbered
  stringstyle=\color{red},     % string literal style
  tabsize=4,                       % sets default tabsize to 2 spaces
  title=\lstname                   % show the filename of files included with \lstinputlisting; also try caption instead of title
}

\newtheoremstyle{custom}
{\topsep}   % Space above
{\topsep}   % Space below
{\slshape}  % Body font
{0pt}       % Indent amount (empty value is the same as 0pt)
{\bfseries} % Theorem head font
{}          % Punctuation after theorem head
{5pt plus 1pt minus 1pt} % Space after theorem head
{\thmname{#1}\thmnumber{ #2}\thmnote{ {\it #3}}\\} % Theorem head spec

\theoremstyle{custom}
\newtheorem{mydef}{Definición}
\numberwithin{mydef}{chapter}

\addto\captionsspanish{
\renewcommand{\abstractname}{Resumen}
\renewcommand{\tablename}{\bf Tabla}
\renewcommand{\partname}{Parte}
\renewcommand{\appendixname}{Anexo}
\renewcommand{\figurename}{\bf Figura}
\renewcommand{\chaptername}{Capítulo}
\renewcommand{\contentsname}{Índice}
\renewcommand{\listtablename}{Lista de Tablas}
\renewcommand{\listfigurename}{Lista de Figuras}
\renewcommand{\bibname}{Bibliografía}
\floatname{algorithm}{Algoritmo}
\renewcommand{\algorithmicrequire}{\textbf{Entrada:}}
\renewcommand{\algorithmicensure}{\textbf{Salida:}}
}

\def\myauthor{Luciano Arnaldo Romero Calla}
\def\mytitle{Template LaTeX tesis Título Profesional CSUNSA}

\title
{
	\includegraphics[width=0.32\textwidth]{figures/unsa}\\\vspace{1cm}
	\large UNIVERSIDAD NACIONAL DE SAN AGUSTÍN\\ 
	Escuela Profesional de Ciencia de la Computación\\\vspace{3cm}
	\Large\bf\mytitle\\\vspace{2cm}
}

\author{\Large\sc\myauthor}

\date{\vspace{3cm}\sc Arequipa\\2017}


\def\figura#1#2#3
{
\begin{figure}[!h]
\centering\includegraphics[width = #2\textwidth]{figures/#1}
\caption{\em #3}\label{#1}
\end{figure}
}

% size, figures array, label, caption
\def\images#1#2#3#4
{
	\begin{figure}[!htb]
		\centering
		\foreach \imagepath/\imagename/\imagecaption in #2
		{
			\subcaptionbox{\label{\imagepath/\imagename} \em \imagecaption}
			{\includegraphics[width=#1\textwidth]{\imagepath/\imagename}}
		}
		\caption{\em #4}\label{#3}
	\end{figure}
}

\def\myimage#1#2%
{%
	\begin{figure}[!htb]
		\centering
		\input{#1}
		\caption{\em #2}\label{#1}
	\end{figure}
}

\newcommand{\tabla}[3]
{
\begin{table}[!h]
\centering#3
\caption{\em #2}\label{#1}
\end{table}
}

\def\refcap#1{\hyperref[#1]{Capítulo \ref{#1}}}
\def\refsec#1{\hyperref[#1]{Sección \ref{#1}}}
\def\reffig#1{\hyperref[#1]{Figura \ref{#1}}}
\def\reftab#1{\hyperref[#1]{Tabla \ref{#1}}}

%%%%%%%%%%%%%%%%%%%%%%%%%%%%%%%%%%%%%%%%%%%%%%%%%%%%%%%%%%%%%%%%%%%%%%%%%%%%%%%%%%%%%%%%%%%%%%%%%%%%%%%%%%%%%%%%%%%%%%%%%%%%%%%%%%

\begin{document}

\maketitle

\begin{titlepage}
\begin{center}
{\Large\sc\myauthor}\\\vspace{2cm}
{\Large\mytitle}\\\vspace{3cm}
\end{center}
\null\hfill
\begin{minipage}{0.48\textwidth}
\large
Tesis presentada para optar el Título Profesional de
Licenciado en Ciencia de la Computación.
\end{minipage}

\bigskip\vspace{2cm}
\large Aprobada en Mayo del 2017.\\\vspace{2cm}

\null\hfill
\begin{minipage}{0.52\textwidth}
\large
\begin{itemize}
\item[] \textbf{Miembros del Jurado:}
\item[\bf Asesor:] Prof. Cristian José López Del Álamo
\item[\bf Presidente:] Prof. Wilber Ramos Lovón
\item[\bf Secretario:] Prof. Juan Carlos Gutiérrez Cáceres
\end{itemize}
\end{minipage}

\bigskip\vspace{4cm}
\begin{center}
\large\sc Arequipa\\2017
\end{center}
\end{titlepage}

\begin{titlepage}
\null\vfill\null\hfill\em Este trabajo está dedicado a mis padres.
\end{titlepage}

\setlength{\parskip}{1em}
\begin{abstract}
\lipsum
\end{abstract}


%\chapter*{Agradecimiento}

\lipsum


\setlength{\parskip}{0em}
\tableofcontents
\listoffigures
\listoftables

\setlength{\parskip}{1em}
% Introduction and Background
\chapter{Introducción}

\lipsum 
\figura{csunsa}{0.6}{CS UNSA\cite{FuentesPerez2017}}

\section{Contribuciones}

\lipsum

\section{Estructura de la tesis}

\lipsum

\chapter{Estado del Arte}\label{background}

\lipsum
			% this chapter should only present fundamental knowledge ...
									% ... necessary in order to understand the problem (but not necessarily the solution)
\chapter{Problema}\label{problem}

\lipsum
			% problem description and problem statement

% Main body – Describing your work
\chapter{Marco Teórico}\label{theory}

\lipsum
				% (advanced preliminaries and/or theoretical part of solution).
\chapter{Propuesta}\label{approach}

\lipsum
			% Description of approach and method(s) to solve the problem
\chapter{Experimentos y Resultados}\label{result}

\lipsum

				% result analysis

% Wrap-up
%\chapter{Trabajos Relacionados}

\lipsum
		% It should demonstrate the principal differences and
									% similarities with respect to 
									% (1) the details of the problem, 
									% (2) the approach, 
									% (3) the results, and possibly 
									% (4) the methodology
\chapter{Conclusiones y trabajos futuros}\label{conclusions}
\lipsum
		% presents conclusions and a final analysis, and possibly some afterwords

% this Structure is from.- Thesis Projects: A Guide for Students in Computer Science and Information Systems
% Mikael Berndtsson, Jörgen Hansson, B. Olsson, Björn Lundell, Second Edition Springer-Verlag London Limited 2008

\bibliographystyle{IEEEtran}
\bibliography{tesis}

\appendix
\chapter{Resultados}

\lipsum


\end{document}

